\section{Related work}
\label{sec:related}

\emph{The computational line.} The paper we comment on \cite{refpaper} inherits
its three-block separation and its max-det SDP from \cite{tanaka2018di}. SDPs in
this literature compute a \emph{value} on a fixed source --- the nonanticipative
rate-distortion function of a partially observable Gauss--Markov process is SDP
computable when strictly feasible \cite{stavrou2019indirect} --- and we claim no
counterpart: the one certificate we built by boxing eigenvalues is retracted in
Section~\ref{sec:waterfill} because it violated its own validity gate, so every
frontier number below is constructive unless a lower bound is quoted. The joint
design paper \cite{tanaka2015joint} removes the fixed $F$ instead: it prices
information in the objective rather than bounding the cost, and it closes by
asking whether the linearity of the sensor policy is restrictive. Sections
\ref{sec:repro}--\ref{sec:nofactor} answer part of that question for one plant and
one cone, which is not the same as resolving it. Two of its moving parts are ours
and we say so plainly: its Step~3 realizability test
$C_tV_t^{-1}C_t^{\top}=P_{t|t}^{-1}-P_{t|t-1}^{-1}$ is the same equality-type
reachability condition we use, so we claim no new construction there; and where it
uses monotonicity of the determinant, the object is one free PSD variable held
under a Loewner upper bound, which is a different statement from the
$\det\Gamma$ monotonicity of Section~\ref{sec:waterfill}: that one is
\emph{false} (a closed-form rank-one counterexample plus $115$ reachable in-cone
counterexamples), and the frontier-wide bound that needed it is withdrawn, not
``unproved''. Distinguishing the two is what keeps a referee from reading this
note as a contradiction of \cite{tanaka2018di}.

\emph{Where the restriction sits.} \cite{cuvelier2021side} lets the decoder hold a
state subset for free, which is a task-restricted virtual sensor with a degenerate
$F$, and is the closest prior for the framing itself; what this note adds over it
is the cone language and the length of the infeasible interval, not the act of
restricting sensing. \cite{sabag2021reducing} already plots a frontier for a
restricted observation matrix and reads its large-cost constant off the unstable
modes left unobservable (Remark~\ref{rmk:dual}). A second selection line chooses
\emph{which} sensor to trust rather than \emph{how much} to say:
\cite{li2023learning,li2026oss} pick one of $N$ sensors at each step or trial, where
every sensor observes the full state with isotropic noise of unknown level, and they
budget trials rather than bits. Their reduction to ``use the smallest covariance all
the time'' is stated in the introduction, not proved as a numbered result, and it is
monotonicity of the expected loss \emph{in the noise level} with the sensing map
already fixed --- their loss carries the covariance through
$\sum_k\operatorname{tr}[S(k)]$ linearly, and covariances comparable in the cone are
in their model just ordered scalars. That is a different statement from the
design-dependent monotonicity this note withdrew, and neither paper prices a rank
choice: their guarantees are a closed-form expected loss and almost-sure value
convergence at rate $t^{-1/2}$, with no frontier and no converse.
\cite{fox2016a,fox2016b} constrain
the controller rather than the sensor and obtain spectral thresholding and an order
phase transition there; our rank loss is on the sensing side, and the switch-on
rates of Section~\ref{sec:waterfill} are stated as measured behaviour of an
upper-bounding family, so none of their phase-transition language transfers to this
note. \cite{tzoumas2020co} pays for sensors rather than for bits and its hard half
is combinatorial---NP-hard already for one step and unit costs---so it
characterises approximation ratios for a selection problem whose trade-off shape is
deliberately left uncharacterised; that shape is what Sections
\ref{sec:infeasible}--\ref{sec:nofactor} measure.

\emph{Rate side.} The stabilizing constant
$\sum_{|\lambda_i(A)|>1}\log_2|\lambda_i(A)|$ that appears in
Section~\ref{sec:repro} and in Remark~\ref{rmk:dual} is a classical quantity: the
finite-data-rate stabilizability bound of \cite{nair2004stabilizability}, read
through the anytime-capacity lens of \cite{sahai2006anytime} and printed against an
LQG frontier in \cite{sabag2021reducing}. It is cited here, not claimed. The
rate-side converse closest to our formulation is the rate-cost function of
\cite{kostina2019rate},
which minimises directed information subject to a cost cap rather than the other
way round; in the indirect NRDF of \cite{stavrou2019indirect} the observation
matrix $C_t$ in $z_t=C_tx_t+n_t$ is non-random with $m\le p$, and the distortion
priced is an unweighted trace above its own floor, so no converse there is a
function of a chosen sensing subspace. Synthesising an observation process is
claimed for the \emph{test channel} --- \cite{stavrou2016estimation} designs the
encoder--channel--decoder triple and states that the sensor must be designed,
but its application is a fully observed Gauss--Markov process with estimation MSE,
pricing no control action; the threshold it derives is a per-component on/off rule
indexed by time and by the eigendirections of the estimator's own error covariance,
so the sensor remains an implementation of a chosen test channel rather than a
design variable priced against a control cost. That floor is not the asymptote studied here, but a
published converse returns it: \cite{li2026sensing} bound, for arbitrary
observation laws including rank-deficient and nonlinear ones, the directed
information rate from the unstable sub-block below by the sum above (their
statement sums over $|\lambda_i(A)|\ge1$), which is the level our designs
approach. Their premise is mean-square observability, so they price nothing on the
cost side and are silent where the filter of Section~\ref{sec:infeasible}
diverges. We offer no closed-form frontier, and the reason is on the record: the
closed-form dynamic reverse water-filling of \cite{tatikonda2004stochastic} for
multidimensional Gauss--Markov sources is false in general:
\cite{stavrou2017revisited} give a two-dimensional counterexample and identify the
equal-distortion assignment as the defect, valid only when the source is i.i.d.;
the corrected object is an SDP value \cite{tanaka2017srd}. Two further bounds are
orthogonal to the cost-side floor: \cite{li2026lower} on multi-encoder networks
without feedback, \cite{atay2026rate} on nonlinear plants with a constant gap. The
closest paper that plots cost against a communication \emph{rate} and proves an
ordering is \cite{shi2012scheduling}, but its rate is the duty cycle
$\limsup_N\frac1N\sum_k I_k$ of one sensor that sees the whole state, its ``optimal''
schedule is optimal among open-loop periodic ones, and its statements run through
traces and L\"owner orderings only --- no determinant appears, and no converse at
fixed rate. The ordering we invert in Section~\ref{sec:nofactor} is a statement about
\emph{different designs} at a fixed rate, so that paper neither corroborates nor
contradicts it. The rate axis plots directed information, which lower-bounds the bits per step any
prefix-free code needs \cite{tanaka2016prefix}; every rate here is read that way.
Finally, the $\Phi$ of \cite{pacelli2026limits} is that paper's own
fixed-point operator; $\Phi_0$ in Section~\ref{sec:infeasible} and $\Phi$ in
Section~\ref{sec:waterfill} are local to this note.
